% −\textit{− coding: utf−8 −}−
% Section 1: Core Principles
% Comprehensive Resource Guide for Fate's Edge

\chapter{Core Principles}łabel{chap:coreₚrinciples}

\section{Safety Tools}łabel{sec:safetyₜools}

Fate's Edge explores mature themes−−−−−−conflict, loss, horror, betrayal. The table should always feel safe. Use these tools to keep the game enjoyable for everyone.

\begin{tcolorbox}[title=Safety Tools: Lines, Veils & X−Card, colback=red!5!white, colframe=red!75!black, sharp corners]
\textbf{X−Card}∈dex{Safety Tools!X−Card} Any player (including the GM) can draw an "X" on a card or simply say "X" when a moment makes them uncomfortable. When the X−Card is shown, the described content is removed from the fiction−−−−−−no questions, no explanations, no judgments. The scene is retconned or redirected immediately.

\textbf{Lines}∈dex{Safety Tools!Lines} A hard boundary: content that will not appear in the game at all. Examples: "No sexual violence," "No harm to children," "No body horror." Establish during Session Zero.

\textbf{Veils}∈dex{Safety Tools!Veils} A soft boundary: content that can happen "off screen" or fade to black. Example: "A romantic scene can happen, but fade to black before it becomes explicit."  
\end{tcolorbox}

Always keep an X−Card visible. When used, pause, rewind if needed, and move on. Do not interrogate the player. The GM may also use the X−Card for their own comfort.

\section{The Central Question}łabel{sec:central_question}

At its heart, Fate's Edge asks:∈dex{Core Philosophy!Central Question}
\begin{quote}
\textbf{What are you willing to risk, and what are you willing to surrender, to reshape the world around you?}
\end{quote}
This question is both philosophical and mechanical. Players gamble with fate every time they act, and the consequences−−−−−−good or ill−−−−−−become the foundation of their legend.

\section{Core Principles}łabel{sec:coreₚrinciplesₗist}

The game is built on four key ideas:∈dex{Core Philosophy!Principles}

\begin{description}
ℹtem[Narrative First]∈dex{Core Philosophy!Narrative First} Mechanics serve the story. Rules reward descriptive play and creative problem−solving.
ℹtem[Risk Creates Drama]∈dex{Core Philosophy!Risk Creates Drama} Every roll carries tension. Even success may come at a cost.
ℹtem[Meaningful Growth]∈dex{Core Philosophy!Meaningful Growth} Experience is a currency of choice. Invest in yourself or in your influence on the world.
ℹtem[Consequences Matter]∈dex{Core Philosophy!Consequences Matter} No action is free. Every choice changes the fiction.
\end{description}

\section{Key Concepts}łabel{sec:key_concepts}

\subsection{Narrative Time}łabel{subsec:narrativeₜime}∈dex{Narrative Time}

Time in Fate's Edge is measured by story weight, not by timers.∈dex{Narrative Time!Scales} Use these four scales:

\begin{longtable}{p{0.2łinewidth} p{0.7łinewidth}}
∩tion{Narrative Time Scales}łabel{tab:narrativeₜime}\\
⊤rule
\textbf{Scale} & \textbf{Description} \\
∣rule
\endfirsthead
μlticolumn{2}{c}%
{{\bfseries \tablename\ θble{} −− continued}} \\
⊤rule
\textbf{Scale} & \textbf{Description} \\
∣rule
\endhead
⊥tomrule
\endfoot
A Moment & A heartbeat, a glance, a single strike or word.∈dex{Narrative Time!Moment} \\
Some Time & A few minutes: a skirmish, a careful lockpick, a short negotiation.∈dex{Narrative Time!Some Time} \\
Significant Time & Hours: travel, a ritual, a siege.∈dex{Narrative Time!Significant Time} \\
Days & Large−scale endeavors: marches, training, recovery.∈dex{Narrative Time!Days} \\
\end{longtable}

\subsection{Story Beats (SB)}łabel{subsec:story_beats}∈dex{Story Beats}

Whenever a player rolls dice, each result of \textbf{1} generates one Story Beat (SB) for the GM.∈dex{Story Beats!Generation} Re−rolling a 1 does \textit{not} erase its SB; if the re−rolled die also shows 1, it generates additional SB.

SB are not penalties−−−−−−they are narrative fuel. The GM spends SB to introduce twists:

\begin{itemize}
ℹtem \textbf{Escalation}∈dex{Story Beats!Escalation} −− drawing more enemies, raising the stakes.
ℹtem \textbf{Exhaustion}∈dex{Story Beats!Exhaustion} −− draining time, resources, or positioning.
ℹtem \textbf{Exposure}∈dex{Story Beats!Exposure} −− revealing hidden actions, alerting foes.
ℹtem \textbf{Collateral}∈dex{Story Beats!Collateral} −− harm or danger spilling over onto allies, innocents, or surroundings.
\end{itemize}

For a full SB Spend Menu, see \ref{sec:sbₛpendₘenu}.

\subsection{Affinity}łabel{subsec:affinity}∈dex{Affinity}

Each people (culture or heritage) provides an \textbf{Affinity}: a narrative edge or metaphysical bond.∈dex{Character Creation!Affinity} Affinities make certain actions more reliable or grant unique permissions, weaving identity into mechanics. Examples include the Aeler stone−sense or the Aelaerem hearth−luck.∈dex{Races!Affinities}

\subsection{Step 6 – Cultural Affinity}

\paragraph{Note on Stacking.}
If your ancestral gift (from your people) and your cultural affinity provide 
the same benefit (e.g., both improve Position), they do not stack. You may 
use one per scene. Choose which applies when you need it.
\subsection{Prestige Abilities}łabel{subsec:prestigeₐbilities}∈dex{Prestige Abilities}

High−level talents unlocked by mastering cultural arts, Patron bonds, or philosophical paths.∈dex{Advancement!Prestige Abilities} These are narrative milestones as much as mechanical upgrades.

\subsection{On−Screen vs. Off−Screen}łabel{subsec:onscreenₒffscreen}∈dex{Resources!On−Screen vs. Off−Screen}

Fate's Edge distinguishes between resources you see at the table and those that shape the world behind the scenes:

\begin{description}
ℹtem[On−Screen Resources]∈dex{Resources!On−Screen} are companions, hirelings, or allies who stand beside you in danger. They add dice pools and flavor, but they can falter, be captured, or die.
ℹtem[Off−Screen Resources]∈dex{Resources!Off−Screen} are assets, titles, networks, or safe houses. They never swing a blade in combat, but they shape the story between sessions, turning XP into narrative leverage.
\end{description}

\section{Design Philosophy}łabel{sec:designₚhilosophy}

\subsection{Core Principles (Extended)}łabel{subsec:design_core}

\begin{enumerate}
    ℹtem \textbf{Narrative Primacy}∈dex{Design Philosophy!Narrative Primacy}: Mechanics serve story, not replace it.
    ℹtem \textbf{Risk as Drama}∈dex{Design Philosophy!Risk as Drama}: Every roll carries potential for triumph and complication.
    ℹtem \textbf{Meaningful Growth}∈dex{Design Philosophy!Meaningful Growth}: XP investment creates lasting character and world change.
    ℹtem \textbf{Consequence Weight}∈dex{Design Philosophy!Consequence Weight}: Choices ripple outward; nothing is free.
    ℹtem \textbf{Fail Forward}∈dex{Design Philosophy!Fail Forward}: Misses fuel Boons; 1s become Story Beats that drive the fiction onward.
\end{enumerate}

\subsection{Mechanical Constraints}łabel{subsec:mechanical_constraints}∈dex{Design Philosophy!Mechanical Constraints}

\begin{longtable}{p{0.3łinewidth} p{0.65łinewidth}}
∩tion{Core Mechanical Limits}łabel{tab:mechₗimits}\\
⊤rule
\textbf{Rule} & \textbf{Limit} \\
∣rule
\endfirsthead
μlticolumn{2}{c}%
{{\bfseries \tablename\ θble{} −− continued}} \\
⊤rule
\textbf{Rule} & \textbf{Limit} \\
∣rule
\endhead
⊥tomrule
\endfoot
Assist Maximum & \(+3\) dice total per roll from all helpers.∈dex{Mechanics!Assist Maximum} \\
Boon Maximum & Hold 5 Boons; trim to 2 at scene end.∈dex{Boons!Maximum} \\
Boon−to−XP Conversion & 2 Boons \(→\) 1 XP, once per session, max 2 XP per session via conversion.∈dex{Boons!Conversion} \\
Follower Initiative & 1 Follower Action per scene, party−wide.∈dex{Combat!Follower Initiative} \\
Over−Stack & 2+ structural advantages = start a relevant timer \(+1\) OR GM banks \(+1\) SB.∈dex{Combat!Over−Stack} \\
Position States & Dominant | Controlled | Desperate (affects success/failure texture).∈dex{Combat!Position States} \\
Dice Pool Cap & Maximum 10d10 per roll; extras become auto−successes.∈dex{Mechanics!Dice Cap} \\
\end{longtable}

\subsection{The Golden Rule}łabel{subsec:goldenᵣule}∈dex{The Golden Rule}

\begin{quote}
\textit{When in doubt, make the choice that serves the story. If you find yourself asking "What's the optimal move?" try reframing it as "What would be most interesting for the story?"}
\end{quote}

The dice do not determine success or failure−−−−−−they help decide which path you take to get there.

\subsection{Flavor is Free}łabel{subsec:flavorᵢs_free}

Players may add descriptive details, cultural elements, and atmospheric touches without spending resources or requiring rolls. Flavor enriches the narrative without changing mechanical outcomes. The GM should encourage this and reciprocate.∈dex{Game Mastering!Flavor is Free}

% End of Section 1